\appendix

\section{Artifact Appendix}

\subsection{Abstract}
%
In this work, we present \sysname{}, a processing-in-memory
acceleration solution for post-transformer large language models.
%
Our artifact includes a full-system simulator for \sysname{} as well as
accuracy evaluation code.
%
To ensure reproducibility and streamline the experimental workflow,
we made two key engineering efforts:
%
(1) The entire codebase relies solely on the modern project manager
\texttt{uv}, allowing evaluators to install exactly the same versions
of dependencies and build tools with a one-line command; and
%
(2) we provide only two scripts--one to run all experiments and
another to generate all figures from the results--to reduce the burden
on evaluators.
%
These two engineering efforts collectively enable a streamlined
experimental workflow and maximize reproducibility.
%

%
Beyond the artifact evaluation, we also document the challenges we
encountered while preparing
the code and a breif API reference to assist those who wish to extend
our codebase.
%
We believe these resources will help those preparing artifact
evaluations and aid future research efforts.
%

\subsection{Artifact check-list}

\begin{itemize}[leftmargin=*]
  \item \textbf{Program:} uv
  \item \textbf{Compilation:} gcc
  \item \textbf{Model:} Our code automatically downloads all required
    model weights from Hugging Face: RetNet, GLA, HGRN2, Mamba2,
    Zamba2, OPT, LLaMA
  \item \textbf{Data sets:} Our code automatically downloads all required
    datasets from Hugging Face: Wikitext2, Piqa, Lambada,
    Hellaswag, Arc-easy, Arc-challenge, Winogrande
  \item \textbf{Run-time environment:} CUDA, x86\_64
  \item \textbf{Hardware:} An NVIDIA GPU based on the Ampere
    architecture or newer, equipped with at least 24GB of memory
  \item \textbf{Metrics:} Perplexity, Accuracy, Throughput
  \item \textbf{Output:} PDF files for the figures, a CSV file for the table
  \item \textbf{How much disk space required (approximately)?:} The
    installation occupies about 100GB in the user's "\$HOME" directory,
    primarily due to models and datasets stored under
    "\$HOME/.cache/huggingface".
  \item \textbf{How much time is needed to prepare workflow
    (approximately)?:} 10 minutes
  \item \textbf{How much time is needed to complete experiments
    (approximately)?:} 15 hours
  \item \textbf{Publicly available?:}
    \bluetext{\url{https://github.com/casys-kaist/pimba}}
  \item \textbf{Archived (provide DOI)?:}
    \bluetext{\url{https://doi.org/10.5281/zenodo.16946084}}
\end{itemize}

%%%%%%%%%%%%%%%%%%%%%%%%%%%%%%%%%%%%%%%%%%%%%%%%%%%%%%%%%%%%%%%%%%%%
\subsection{Description}
%
\niparagraph{How to access.}
%
Clone our public repository from github:
%
\begin{verbatim}
  $ git clone https://github.com/casys-kaist/pimba
\end{verbatim}
%

%
\niparagraph{Hardware dependencies.}
%
Evaluators require an NVIDIA GPU based on
the Ampere architecture or newer, with at least 24GB of memory
(e.g., RTX 3090, RTX 4090, RTX 5090, A6000, A100, etc.).
%

%
\niparagraph{Software dependencies.}
%
Evaluators only need \texttt{gcc} and \texttt{uv}; all other required
python packages and build tools are installed through \texttt{uv}.
%
\texttt{uv} itself can be installed using the system package manager
or via a one-line command provided in the official documentation~\cite{uv}.
%

%
\niparagraph{Models and data sets.}
%
During the experiments, our code automatically downloads all required
models and datasets.
%

\subsection{Installation}
%
With the following commands, \texttt{uv} automatically downloads the required
dependencies and build tools (e.g., cmake, ninja) in exactly the same
versions we used and compiles our project.
%
\begin{verbatim}
  $ uv sync
  $ uv run cmake --preset release
  $ uv run cmake --build build
\end{verbatim}

\subsection{Experiment workflow}
%
Evaluators can run all experiments with a single command:
%
\begin{verbatim}
  $ uv run python scripts/run.py
\end{verbatim}
%
This command generates two files, "accuracy\_result.yaml" and
"performance\_result.yaml", under the "res/" directory, which are
then used to reproduce the figures and the table.

\subsection{Evaluation and expected results}
%
Using the result files, evaluators can simply reproduce the figures
and the table in the paper by executing the following command:
%
\begin{verbatim}
  $ uv run python scripts/draw.py
\end{verbatim}
%
This process generates PDF files for the figures and a CSV file for
the table in the "summary/" directory.
%
Note that the PDF files do not include axis labels or legends, as we
add them using external drawing tools.
%
However, all scales and colors precisely match those in the paper,
ensuring no issues arise when
comparing values.
%

%
We also observed that accuracy may slightly vary with the GPU
platform and CUDA version used for evaluation, especially in
Table~\ref{tab:accuracy}.
%
However, the key trend we aim to convey in this paper remains
consistent: quantization with \texttt{MX8} has minimal impact on accuracy.
%
